% Options for packages loaded elsewhere
\PassOptionsToPackage{unicode}{hyperref}
\PassOptionsToPackage{hyphens}{url}
%
\documentclass[
]{article}
\usepackage{amsmath,amssymb}
\usepackage{iftex}
\ifPDFTeX
  \usepackage[T1]{fontenc}
  \usepackage[utf8]{inputenc}
  \usepackage{textcomp} % provide euro and other symbols
\else % if luatex or xetex
  \usepackage{unicode-math} % this also loads fontspec
  \defaultfontfeatures{Scale=MatchLowercase}
  \defaultfontfeatures[\rmfamily]{Ligatures=TeX,Scale=1}
\fi
\usepackage{lmodern}
\ifPDFTeX\else
  % xetex/luatex font selection
\fi
% Use upquote if available, for straight quotes in verbatim environments
\IfFileExists{upquote.sty}{\usepackage{upquote}}{}
\IfFileExists{microtype.sty}{% use microtype if available
  \usepackage[]{microtype}
  \UseMicrotypeSet[protrusion]{basicmath} % disable protrusion for tt fonts
}{}
\makeatletter
\@ifundefined{KOMAClassName}{% if non-KOMA class
  \IfFileExists{parskip.sty}{%
    \usepackage{parskip}
  }{% else
    \setlength{\parindent}{0pt}
    \setlength{\parskip}{6pt plus 2pt minus 1pt}}
}{% if KOMA class
  \KOMAoptions{parskip=half}}
\makeatother
\usepackage{xcolor}
\usepackage[margin=1in]{geometry}
\usepackage{color}
\usepackage{fancyvrb}
\newcommand{\VerbBar}{|}
\newcommand{\VERB}{\Verb[commandchars=\\\{\}]}
\DefineVerbatimEnvironment{Highlighting}{Verbatim}{commandchars=\\\{\}}
% Add ',fontsize=\small' for more characters per line
\usepackage{framed}
\definecolor{shadecolor}{RGB}{248,248,248}
\newenvironment{Shaded}{\begin{snugshade}}{\end{snugshade}}
\newcommand{\AlertTok}[1]{\textcolor[rgb]{0.94,0.16,0.16}{#1}}
\newcommand{\AnnotationTok}[1]{\textcolor[rgb]{0.56,0.35,0.01}{\textbf{\textit{#1}}}}
\newcommand{\AttributeTok}[1]{\textcolor[rgb]{0.13,0.29,0.53}{#1}}
\newcommand{\BaseNTok}[1]{\textcolor[rgb]{0.00,0.00,0.81}{#1}}
\newcommand{\BuiltInTok}[1]{#1}
\newcommand{\CharTok}[1]{\textcolor[rgb]{0.31,0.60,0.02}{#1}}
\newcommand{\CommentTok}[1]{\textcolor[rgb]{0.56,0.35,0.01}{\textit{#1}}}
\newcommand{\CommentVarTok}[1]{\textcolor[rgb]{0.56,0.35,0.01}{\textbf{\textit{#1}}}}
\newcommand{\ConstantTok}[1]{\textcolor[rgb]{0.56,0.35,0.01}{#1}}
\newcommand{\ControlFlowTok}[1]{\textcolor[rgb]{0.13,0.29,0.53}{\textbf{#1}}}
\newcommand{\DataTypeTok}[1]{\textcolor[rgb]{0.13,0.29,0.53}{#1}}
\newcommand{\DecValTok}[1]{\textcolor[rgb]{0.00,0.00,0.81}{#1}}
\newcommand{\DocumentationTok}[1]{\textcolor[rgb]{0.56,0.35,0.01}{\textbf{\textit{#1}}}}
\newcommand{\ErrorTok}[1]{\textcolor[rgb]{0.64,0.00,0.00}{\textbf{#1}}}
\newcommand{\ExtensionTok}[1]{#1}
\newcommand{\FloatTok}[1]{\textcolor[rgb]{0.00,0.00,0.81}{#1}}
\newcommand{\FunctionTok}[1]{\textcolor[rgb]{0.13,0.29,0.53}{\textbf{#1}}}
\newcommand{\ImportTok}[1]{#1}
\newcommand{\InformationTok}[1]{\textcolor[rgb]{0.56,0.35,0.01}{\textbf{\textit{#1}}}}
\newcommand{\KeywordTok}[1]{\textcolor[rgb]{0.13,0.29,0.53}{\textbf{#1}}}
\newcommand{\NormalTok}[1]{#1}
\newcommand{\OperatorTok}[1]{\textcolor[rgb]{0.81,0.36,0.00}{\textbf{#1}}}
\newcommand{\OtherTok}[1]{\textcolor[rgb]{0.56,0.35,0.01}{#1}}
\newcommand{\PreprocessorTok}[1]{\textcolor[rgb]{0.56,0.35,0.01}{\textit{#1}}}
\newcommand{\RegionMarkerTok}[1]{#1}
\newcommand{\SpecialCharTok}[1]{\textcolor[rgb]{0.81,0.36,0.00}{\textbf{#1}}}
\newcommand{\SpecialStringTok}[1]{\textcolor[rgb]{0.31,0.60,0.02}{#1}}
\newcommand{\StringTok}[1]{\textcolor[rgb]{0.31,0.60,0.02}{#1}}
\newcommand{\VariableTok}[1]{\textcolor[rgb]{0.00,0.00,0.00}{#1}}
\newcommand{\VerbatimStringTok}[1]{\textcolor[rgb]{0.31,0.60,0.02}{#1}}
\newcommand{\WarningTok}[1]{\textcolor[rgb]{0.56,0.35,0.01}{\textbf{\textit{#1}}}}
\usepackage{graphicx}
\makeatletter
\newsavebox\pandoc@box
\newcommand*\pandocbounded[1]{% scales image to fit in text height/width
  \sbox\pandoc@box{#1}%
  \Gscale@div\@tempa{\textheight}{\dimexpr\ht\pandoc@box+\dp\pandoc@box\relax}%
  \Gscale@div\@tempb{\linewidth}{\wd\pandoc@box}%
  \ifdim\@tempb\p@<\@tempa\p@\let\@tempa\@tempb\fi% select the smaller of both
  \ifdim\@tempa\p@<\p@\scalebox{\@tempa}{\usebox\pandoc@box}%
  \else\usebox{\pandoc@box}%
  \fi%
}
% Set default figure placement to htbp
\def\fps@figure{htbp}
\makeatother
\setlength{\emergencystretch}{3em} % prevent overfull lines
\providecommand{\tightlist}{%
  \setlength{\itemsep}{0pt}\setlength{\parskip}{0pt}}
\setcounter{secnumdepth}{-\maxdimen} % remove section numbering
\usepackage{bookmark}
\IfFileExists{xurl.sty}{\usepackage{xurl}}{} % add URL line breaks if available
\urlstyle{same}
\hypersetup{
  pdftitle={Identifying Correlated Methylation Units (CMUs)},
  pdfauthor={Aarti Jajoo, Owen Hirchi \& Neil Hanchard},
  hidelinks,
  pdfcreator={LaTeX via pandoc}}

\title{Identifying Correlated Methylation Units (CMUs)}
\author{Aarti Jajoo, Owen Hirchi \& Neil Hanchard}
\date{2025-12-14}

\begin{document}
\maketitle

\subsection{Overview}\label{overview}

Correlated Methylation Units (CMUs) are genomic regions defined by
coordinated DNA methylation patterns across neighboring CpG sites. For a
given 450k or EPIC methylation profile across samples, CMUs are
identified within non-overlapping 250,000 base pair windows across
autosomal chromosomes.

This vignette demonstrates how to:

\begin{enumerate}
\def\labelenumi{\arabic{enumi}.}
\setcounter{enumi}{-1}
\tightlist
\item
  \hyperref[step0-setup]{Setup R environment, dependencies, and load CMU
  functions}
\item
  \hyperref[methylation-input-format]{Prepare Illumina 450k or EPIC DNA
  methylation data}
\item
  \hyperref[step2-indexing]{Generate indexing for efficient CMU
  computation}
\item
  \hyperref[step3-CMU-calling]{Identify CMUs}
\item
  \hyperref[step4-adj_cov]{Incorporate sample subsets and covariates}
\item
  \hyperref[step5-diffCMU]{Perform differential CMUs}
\item
  \hyperref[cmuOutput]{How to use CMU output files}
\end{enumerate}

Toy datasets provided can help in getting the correct format of dataset
to be used.

\begin{center}\rule{0.5\linewidth}{0.5pt}\end{center}

\subsection{Step 0: Setup environment, dependencies, and load CMU
functions}\label{step0-setup}

The CMU pipeline is implemented primarily in R, with Python dependency
used for indexing methylation data.

\begin{Shaded}
\begin{Highlighting}[]
\FunctionTok{library}\NormalTok{(reticulate)}
\FunctionTok{library}\NormalTok{(matrixStats)}
\FunctionTok{library}\NormalTok{(spatstat)}
\FunctionTok{library}\NormalTok{(imager)}
\CommentTok{\#Python dependency installation:}
\NormalTok{reticulate}\SpecialCharTok{::}\FunctionTok{py\_install}\NormalTok{(}\StringTok{"pandas"}\NormalTok{)}
\end{Highlighting}
\end{Shaded}

\begin{quote}
\textbf{Note for macOS users:} The \texttt{imager} package may require
XQuartz. If pacakge installation fails, install XQuartz from
\url{https://www.xquartz.org/} and restart R.
\end{quote}

🛑 Do not use \textasciitilde{} (tilde) in the filepath i.e.~isntead of
\textasciitilde/Data use /Users/ajajoo/Data. Otherwise, python
reticulate related commands give error.

\begin{center}\rule{0.5\linewidth}{0.5pt}\end{center}

\paragraph{Change working directory \& Load CMU
functions}\label{change-working-directory-load-cmu-functions}

Set the working directory to the location where the CMU repository has
been downloaded.

\begin{Shaded}
\begin{Highlighting}[]
\FunctionTok{setwd}\NormalTok{(}\StringTok{"path/to/CMU"}\NormalTok{)}
\end{Highlighting}
\end{Shaded}

\begin{Shaded}
\begin{Highlighting}[]
\FunctionTok{source}\NormalTok{(}\StringTok{"GetIndexing.R"}\NormalTok{)}
\FunctionTok{source}\NormalTok{(}\StringTok{"CUsAcrossGenomeVer05ForPyIndex.R"}\NormalTok{)}
\end{Highlighting}
\end{Shaded}

\begin{center}\rule{0.5\linewidth}{0.5pt}\end{center}

\subsection{Step 1: Methylation input file
formating}\label{methylation-input-format}

CMU analysis requires a \textbf{tab-separated text file} containing DNA
methylation beta values with the following strict structure.

\paragraph{Example (toy data)}\label{example-toy-data}

Example of a beta file:

\begin{Shaded}
\begin{Highlighting}[]
\NormalTok{toyData/EPIC\_Mimic\_Text.tsv}
\end{Highlighting}
\end{Shaded}

\begin{Shaded}
\begin{Highlighting}[]
\NormalTok{"Sample\_1"    "Sample\_2"   }
\NormalTok{"cg26928153"  0.331238528247923   0.869199158390984}
\NormalTok{"cg16269199"  0.0595261559366040  0.934328192146495}
\NormalTok{"cg13869341"  0.354899412253872   0.249734462632203}
\NormalTok{"cg24669183"  0.912160335108638   0.771520412294194}
\end{Highlighting}
\end{Shaded}

Notes:

\begin{itemize}
\tightlist
\item
  \textbf{CpG IDs are quoted}
\item
  \textbf{CpG column does not have a column name in header}
\item
  \textbf{Sample names are quoted}
\item
  \textbf{Values are numeric beta values}
\end{itemize}

Example of a probe file: File location (toy data):

\begin{Shaded}
\begin{Highlighting}[]
\NormalTok{toyData/probe\_Epic.txt}
\end{Highlighting}
\end{Shaded}

\begin{Shaded}
\begin{Highlighting}[]
\NormalTok{"cg26928153"}
\NormalTok{"cg16269199"}
\NormalTok{"cg13869341"}
\NormalTok{"cg24669183"}
\end{Highlighting}
\end{Shaded}

Notes:

\begin{itemize}
\tightlist
\item
  \textbf{One CpG ID per line}

  \begin{itemize}
  \tightlist
  \item
    \textbf{No header}
  \item
    \textbf{CpG IDs must be quoted}
  \item
    \textbf{Order must exactly match the methylation file}
  \end{itemize}
\end{itemize}

\begin{center}\rule{0.5\linewidth}{0.5pt}\end{center}

\begin{quote}
⚠️ \textbf{Input file formatting}\\
If you encounter errors related to file structure, quoting, or probe
ordering,\\
see the detailed formatting requirements and ways to convert your data
in the correct format section\\
\hyperref[detailed-input-data-requirements-and-data-formatting]{Click
here for more details on Input data requirements and data formatting}.
\end{quote}

\begin{center}\rule{0.5\linewidth}{0.5pt}\end{center}

\subsection{Step 2: Generating CMU Folder and Indexing for beta
file}\label{step2-indexing}

\begin{Shaded}
\begin{Highlighting}[]
\CommentTok{\# For EPIC }
\FunctionTok{GenerateDataFormat}\NormalTok{(}
\AttributeTok{filename =} \StringTok{"toyData/EPIC\_Mimic\_Text.tsv"}\NormalTok{,}
\AttributeTok{PathToResult =} \StringTok{"CMUResults/EPIC/"}\NormalTok{,}
\AttributeTok{SortedProbeFilename =} \StringTok{"./PythonIndexing/SortedCpGNamesEPIC.txt"}
\NormalTok{)}
\CommentTok{\# For 450K }
\CommentTok{\# change SortedProbeFilename = "./PythonIndexing/SortedCpGNames450k.txt"}
\end{Highlighting}
\end{Shaded}

This step: - Generates indexing files required for quick reading of beta
file

🛑 DO NOT modify your original beta TSV file after this step.\\
If you do, you must rerun Step 4.

\begin{center}\rule{0.5\linewidth}{0.5pt}\end{center}

\subsection{Step 3: Running CMU detection on Illumina 450k
data}\label{step3-CMU-calling}

\begin{Shaded}
\begin{Highlighting}[]
\FunctionTok{CUsAcrossGenome}\NormalTok{(}
  \AttributeTok{annotationFile =} \StringTok{"Illumina\_probe\_annotation\_Epic.txt"}\NormalTok{,}
  \AttributeTok{PathToFig =} \StringTok{"CMUResults/EPIC/Results/"}\NormalTok{,}
  \AttributeTok{data =} \StringTok{"toyData/EPIC\_Mimic\_Text.tsv"}\NormalTok{,}
  \AttributeTok{probenameFile =} \StringTok{"toyData/probe\_Epic.txt"}\NormalTok{,}
  \AttributeTok{indexFile =} \StringTok{"CMUResults/EPIC/Data/Index.txt"}
\NormalTok{)}
\CommentTok{\# For 450K }
\CommentTok{\# change annotationFile = "Illumina\_probe\_annotation\_450k.txt"}
\end{Highlighting}
\end{Shaded}

This function identifies genome-wide CMUs.

\begin{center}\rule{0.5\linewidth}{0.5pt}\end{center}

\subsection{Step 4: Advanced options: sample subsets and
covariates}\label{step4-adj_cov}

\subsubsection{Group-specific CMU
computation}\label{group-specific-cmu-computation}

You can either prepare a new methylation input following Step 3, or CMUs
can also be computed on a subset of samples by providing a logical
vector (\texttt{GroupIndex}) matching the order of samples in the
methylation matrix.

\begin{Shaded}
\begin{Highlighting}[]
\NormalTok{GroupIndex }\OtherTok{\textless{}{-}}\NormalTok{ SubjectNameInGroup }\SpecialCharTok{\%in\%}\NormalTok{ SubjectNameInBetaTable}
\FunctionTok{CUsAcrossGenome}\NormalTok{(}
    \AttributeTok{annotationFile =} \StringTok{"Illumina\_probe\_annotation\_Epic.txt"}\NormalTok{,}
    \AttributeTok{PathToFig =} \StringTok{"CMUResults/EPIC/Results/"}\NormalTok{,}
    \AttributeTok{data =} \StringTok{"toyData/EPIC\_Mimic\_Text.tsv"}\NormalTok{,}
    \AttributeTok{probenameFile =} \StringTok{"toyData/probe\_Epic.txt"}\NormalTok{,}
    \AttributeTok{indexFile =} \StringTok{"CMUResults/EPIC/Data/Index.txt"}\NormalTok{,}
    \AttributeTok{GroupIndex =}\NormalTok{ GroupIndex}
\NormalTok{  )}
\end{Highlighting}
\end{Shaded}

\subsubsection{Covariate adjustment}\label{covariate-adjustment}

CMUs can be computed after residualizing a choice of user provided
co-variates. The matrix must match the sample order of the input data
and include an intercept column (a column with all 1s).

\begin{Shaded}
\begin{Highlighting}[]
\NormalTok{confounderX }\OtherTok{\textless{}{-}} \FunctionTok{as.matrix}\NormalTok{(}
\NormalTok{  bb[, }\FunctionTok{c}\NormalTok{(}\StringTok{"intercept"}\NormalTok{, }\StringTok{"Gender"}\NormalTok{, }\StringTok{"Age"}\NormalTok{)]}
\NormalTok{)}
\FunctionTok{CUsAcrossGenome}\NormalTok{(}
    \AttributeTok{annotationFile =} \StringTok{"Illumina\_probe\_annotation\_Epic.txt"}\NormalTok{,}
    \AttributeTok{PathToFig =} \StringTok{"CMUResults/EPIC/Results/"}\NormalTok{,}
    \AttributeTok{data =} \StringTok{"toyData/EPIC\_Mimic\_Text.tsv"}\NormalTok{,}
    \AttributeTok{probenameFile =} \StringTok{"toyData/probe\_Epic.txt"}\NormalTok{,}
    \AttributeTok{indexFile =} \StringTok{"CMUResults/EPIC/Data/Index.txt"}\NormalTok{,}
    \AttributeTok{confounderX =}\NormalTok{ confounderX}
\NormalTok{  )}
\end{Highlighting}
\end{Shaded}

\subsection{Step 5: Differential CMUs}\label{step5-diffCMU}

To identify differential CMUs, use DiffCMU.R file. You will have to
provide the following variables.

\begin{Shaded}
\begin{Highlighting}[]
\CommentTok{\# set the following variables }
\NormalTok{data\_file }\OtherTok{\textless{}{-}} \CommentTok{\# methylation tsv file prepared  in Step 1 }
\NormalTok{index\_file }\OtherTok{\textless{}{-}} \CommentTok{\# corresponding index file for  methylation data prepared in Step 2}
\NormalTok{idx\_ctrl\_X  }\OtherTok{\textless{}{-}} \CommentTok{\# index for controls/normal }
\NormalTok{idx\_case\_X  }\OtherTok{\textless{}{-}} \CommentTok{\# index for case/Disease subjects  }
\NormalTok{covX }\OtherTok{\textless{}{-}} \CommentTok{\# covariates to be used, in same order as it is in data\_file. idx\_ctrl\_X and idx\_case\_X would  \# be used to subset this matrix for diseaes and controls. }
\NormalTok{CMU\_file\_for\_ctrl }\OtherTok{\textless{}{-}} \CommentTok{\# path\_to\_CMU\_ctrl\_Output/ContM0.6.txt}
\NormalTok{CMU\_file\_for\_case }\OtherTok{\textless{}{-}} \CommentTok{\# path\_to\_CMU\_case\_Output/ContM0.6.txt}
\NormalTok{pvalFilenameForNormalCMU }\OtherTok{\textless{}{-}} \CommentTok{\# Provide filename with path and extension Rdata}
\NormalTok{pvalFilenameForDiseaseCMU }\OtherTok{\textless{}{-}} \CommentTok{\# Provide filename with path and extebsuib Rdata}
\end{Highlighting}
\end{Shaded}

\subsection{How to use CMU output files}\label{cmuOutput}

CMU calling is performed with three correlation soft thresholds: 0.4,0.6
and 0.8. We recommend 0.6 for general purposes and refer to our paper
for detailed understanding of how lower or higher parameter can impact
the CMU calling. For each parameter their are two files, e.g.~for 0.4 we
have ContM0.6.txt for contiguous CMUs and NonContM0.6.txt for
non-contiguous CMUs.

\begin{Shaded}
\begin{Highlighting}[]
\CommentTok{\# }
\FunctionTok{source}\NormalTok{(}\StringTok{\textquotesingle{}getClusDetails.R\textquotesingle{}}\NormalTok{)}
\NormalTok{CMUs }\OtherTok{\textless{}{-}} \FunctionTok{read.table}\NormalTok{(}\StringTok{\textquotesingle{}pathToCMU/ContM0.6.txt\textquotesingle{}}\NormalTok{,}\AttributeTok{stringsAsFactors =} \ConstantTok{FALSE}\NormalTok{)}
\FunctionTok{library}\NormalTok{(knitr)}
\FunctionTok{kable}\NormalTok{(clusDetailsControl[}\DecValTok{1}\SpecialCharTok{:}\DecValTok{5}\NormalTok{, ], }\AttributeTok{caption =} \StringTok{"Cluster Details"}\NormalTok{)                               V1 V2      V3                                                                                                             V4}
\DecValTok{1}\NormalTok{ ..}\SpecialCharTok{/}\NormalTok{Y123}\SpecialCharTok{/}\NormalTok{Amy}\SpecialCharTok{/}\NormalTok{Results}\SpecialCharTok{/}\NormalTok{ContM0.}\FloatTok{6.}\NormalTok{txt  }\DecValTok{1} \DecValTok{1086777}\NormalTok{ cg24437834;cg15565004;cg19999567;cg19856606;cg12743165;cg04065236;cg04315086;cg09358422;cg24523322;cg12612065;}
\DecValTok{2}\NormalTok{ ..}\SpecialCharTok{/}\NormalTok{Y123}\SpecialCharTok{/}\NormalTok{Amy}\SpecialCharTok{/}\NormalTok{Results}\SpecialCharTok{/}\NormalTok{ContM0.}\FloatTok{6.}\NormalTok{txt  }\DecValTok{1} \DecValTok{1086777}\NormalTok{                                                        cg21786289;cg13897241;cg13176867;cg15105996;cg02565062;}
\DecValTok{3}\NormalTok{ ..}\SpecialCharTok{/}\NormalTok{Y123}\SpecialCharTok{/}\NormalTok{Amy}\SpecialCharTok{/}\NormalTok{Results}\SpecialCharTok{/}\NormalTok{ContM0.}\FloatTok{6.}\NormalTok{txt  }\DecValTok{1} \DecValTok{1086777}\NormalTok{                                                        cg07390924;cg10110857;cg10375192;cg02056921;cg24076968;}
\DecValTok{4}\NormalTok{ ..}\SpecialCharTok{/}\NormalTok{Y123}\SpecialCharTok{/}\NormalTok{Amy}\SpecialCharTok{/}\NormalTok{Results}\SpecialCharTok{/}\NormalTok{ContM0.}\FloatTok{6.}\NormalTok{txt  }\DecValTok{1} \DecValTok{1086777}\NormalTok{                                                        cg20062691;cg03811829;cg11211792;cg04788999;cg16526047;}
\DecValTok{5}\NormalTok{ ..}\SpecialCharTok{/}\NormalTok{Y123}\SpecialCharTok{/}\NormalTok{Amy}\SpecialCharTok{/}\NormalTok{Results}\SpecialCharTok{/}\NormalTok{ContM0.}\FloatTok{6.}\NormalTok{txt  }\DecValTok{1} \DecValTok{1086777}\NormalTok{                                                                   cg22226438;cg08029603;cg22699361;cg03890988;}
\end{Highlighting}
\end{Shaded}

\begin{center}\rule{0.5\linewidth}{0.5pt}\end{center}

\begin{center}\rule{0.5\linewidth}{0.5pt}\end{center}

\subsection{Detailed Input data requirements and data
formatting}\label{detailed-input-data-requirements-and-data-formatting}

\paragraph{Required structure (strict)}\label{required-structure-strict}

\begin{itemize}
\tightlist
\item
  \textbf{Rows:} CpG probe IDs\\
\item
  \textbf{Columns:} Samples\\
\item
  \textbf{Values:} Beta values (0--1)
\end{itemize}

\paragraph{Mandatory formatting rules}\label{mandatory-formatting-rules}

\begin{enumerate}
\def\labelenumi{\arabic{enumi}.}
\tightlist
\item
  \textbf{Tab-separated format}

  \begin{itemize}
  \tightlist
  \item
    The file must be tab-delimited (\texttt{.tsv})
  \item
    Other delimeters are not supported without conversion
  \end{itemize}
\item
  \textbf{First row has one less column}

  \begin{itemize}
  \tightlist
  \item
    The first column contains CpG IDs and therefore has \textbf{no
    column header}
  \item
    So the first row only has sample names which is one column less than
    the remaining files
  \end{itemize}
\item
  \textbf{Quoted identifiers are REQUIRED}

  \begin{itemize}
  \tightlist
  \item
    CpG probe IDs \textbf{must be quoted}
  \item
    Sample names \textbf{must be quoted}
  \item
    Quoting must be consistent across all input files
  \end{itemize}
\item
  \textbf{CpG IDs are read as row names}

  \begin{itemize}
  \tightlist
  \item
    CpG probe IDs appear in the first column
  \end{itemize}
\end{enumerate}

\begin{center}\rule{0.5\linewidth}{0.5pt}\end{center}

\paragraph{Recommended way to get beta file in correct tsv format via
R}\label{recommended-way-to-get-beta-file-in-correct-tsv-format-via-r}

Load your beta file in R, rownames must be CpGIDs and column names must
be sample names.

\begin{Shaded}
\begin{Highlighting}[]
\CommentTok{\# load your methylation file }
\NormalTok{beta\_matrix }\OtherTok{\textless{}{-}} \FunctionTok{read.table}\NormalTok{(}
  \StringTok{"toyData/EPIC\_Mimic\_Text.tsv"}\NormalTok{,}
  \AttributeTok{sep =} \StringTok{"}\SpecialCharTok{\textbackslash{}t}\StringTok{"}\NormalTok{,}
\NormalTok{)}
\end{Highlighting}
\end{Shaded}

Use below command to generate tsv beta file.

\begin{Shaded}
\begin{Highlighting}[]
\FunctionTok{write.table}\NormalTok{(}
\NormalTok{  beta\_matrix,}
  \AttributeTok{file =} \StringTok{"pathtobeta/beta.tsv"}\NormalTok{,}
  \AttributeTok{sep =} \StringTok{"}\SpecialCharTok{\textbackslash{}t}\StringTok{"}
\NormalTok{)}
\end{Highlighting}
\end{Shaded}

\paragraph{Example of converting CSV to TSV
(terminal)}\label{example-of-converting-csv-to-tsv-terminal}

If your methylation data are stored as a comma-separated file
(\texttt{.csv}), convert it to tab-separated format using:

\begin{Shaded}
\begin{Highlighting}[]
\FunctionTok{sed} \StringTok{\textquotesingle{}s/,/\textbackslash{}t/g\textquotesingle{}}\NormalTok{ pathtobeta/beta.csv }\OperatorTok{\textgreater{}}\NormalTok{ pathtobeta/beta.tsv}
\end{Highlighting}
\end{Shaded}

Always verify formatting:

\begin{Shaded}
\begin{Highlighting}[]
\FunctionTok{head}\NormalTok{ pathtobeta/beta.tsv}
\end{Highlighting}
\end{Shaded}

\begin{center}\rule{0.5\linewidth}{0.5pt}\end{center}

\subsubsection{\texorpdfstring{Probe ID file
(\texttt{probe\_Epic.txt})}{Probe ID file (probe\_Epic.txt)}}\label{probe-id-file-probe_epic.txt}

A probe ID file is required and must satisfy the following constraints.

\paragraph{Strict requirements}\label{strict-requirements}

⚠️ \textbf{Critical constraint}

The probe ID file must match the \textbf{exact probe order} of the
methylation file:

\begin{itemize}
\tightlist
\item
  \texttt{probe\_Epic.txt} corresponds to the \textbf{first column} of
  the methylation file
\item
  The methylation file contains \textbf{one additional header row}.
  Therefore, \texttt{probe\_Epic.txt} must contain \textbf{one fewer
  line} compared to beta file.
\end{itemize}

Example \texttt{probe\_Epic.txt}:

\begin{center}\rule{0.5\linewidth}{0.5pt}\end{center}

\#\#\# Generating \texttt{probe\_Epic.txt} from the methylation file
using beta\_matrix (recommended)

\begin{Shaded}
\begin{Highlighting}[]
\FunctionTok{write.table}\NormalTok{(}
\FunctionTok{rownames}\NormalTok{(beta\_matrix),}
\AttributeTok{file =} \StringTok{"toyData/probe\_test.txt"}\NormalTok{,}
\AttributeTok{quote =} \ConstantTok{TRUE}\NormalTok{,}
\AttributeTok{row.names =} \ConstantTok{FALSE}\NormalTok{,}
\AttributeTok{col.names =} \ConstantTok{FALSE}
\NormalTok{)}
\end{Highlighting}
\end{Shaded}

This guarantees: - Exact probe ordering - Quoted CpG IDs - One-to-one
correspondence with the methylation matrix

\begin{center}\rule{0.5\linewidth}{0.5pt}\end{center}

\subsection{Citation}\label{citation}

If you use our CMU software, please cite:

\begin{quote}
\emph{A genome-wide map of correlated DNA methylation developed from
image clustering of multi-tissue methylation data}. (2026).
\end{quote}

\subsection{References}\label{references}

\paragraph{reticulate}\label{reticulate}

Ushey K, Allaire J, Tang Y (2025). \emph{reticulate: Interface to
`Python'}. \url{doi:10.32614/CRAN.package.reticulate}
\url{https://doi.org/10.32614/CRAN.package.reticulate}, R package
version 1.44.0, \url{https://CRAN.R-project.org/package=reticulate}.

\paragraph{matrixStats}\label{matrixstats}

Bengtsson H (2025). \emph{matrixStats: Functions that Apply to Rows and
Columns of Matrices (and to Vectors)}.
\url{doi:10.32614/CRAN.package.matrixStats}
\url{https://doi.org/10.32614/CRAN.package.matrixStats}, R package
version 1.5.0, \url{https://CRAN.R-project.org/package=matrixStats}.

\paragraph{spatstat}\label{spatstat}

Baddeley A, Rubak E, Turner R (2015). \emph{Spatial Point Patterns:
Methodology and Applications with R}. Chapman and Hall/CRC Press,
London. ISBN 9781482210200,
\url{https://www.routledge.com/Spatial-Point-Patterns-Methodology-and-Applications-with-R/Baddeley-Rubak-Turner/p/book/9781482210200/}.
Baddeley A, Turner R, Mateu J, Bevan A (2013). ``Hybrids of Gibbs Point
Process Models and Their Implementation.'' \emph{Journal of Statistical
Software}, \emph{55}(11), 1-43. \url{doi:10.18637/jss.v055.i11}
\url{https://doi.org/10.18637/jss.v055.i11}. Baddeley A, Turner R
(2005). ``spatstat: An R Package for Analyzing Spatial Point Patterns.''
\emph{Journal of Statistical Software}, \emph{12}(6), 1-42.
\url{doi:10.18637/jss.v012.i06}
\url{https://doi.org/10.18637/jss.v012.i06}.

\paragraph{imager}\label{imager}

Barthelme S (2025). \emph{imager: Image Processing Library Based on
`CImg'}. \url{doi:10.32614/CRAN.package.imager}
\url{https://doi.org/10.32614/CRAN.package.imager}, R package version
1.0.5, \url{https://CRAN.R-project.org/package=imager}.

\paragraph{pandas}\label{pandas}

McKinney, W. (2010). Data Structures for Statistical Computing in
Python.

\end{document}
